\documentclass[11pt,a4paper]{article}

% ============================================================
%  QIMMA -- COURS : Fonctions primitives
%  2ème Bac Sciences Mathématiques (A et B)
%  Standard exhaustif : définitions formelles + démonstrations
%  Basé sur templates/qimma-template.tex (charte Qimma)
% ============================================================

\usepackage[french]{babel}
\usepackage[utf8]{inputenc}
\usepackage[T1]{fontenc}
\usepackage{lmodern}
\usepackage[a4paper,margin=2cm,top=2.3cm,bottom=2.6cm]{geometry}
\usepackage{amsmath,amssymb}
\usepackage{xcolor}
\usepackage{graphicx}
\usepackage{tikz}
\usetikzlibrary{shadows,calc}
\usepackage[most]{tcolorbox}
\usepackage{titlesec}
\usepackage{fancyhdr}
\usepackage{enumitem}
\usepackage{pifont}
\usepackage{array}
\usepackage{colortbl}
\usepackage{hyperref}

% ------------------- CHARTE QIMMA -------------------
\definecolor{qteal}{HTML}{0d9488}
\definecolor{qtealdark}{HTML}{134e4a}
\definecolor{qamber}{HTML}{fbbf24}
\definecolor{qambertext}{HTML}{92400e}
\definecolor{ink}{HTML}{1f2937}
\definecolor{inkgray}{HTML}{6b7280}
\definecolor{tealpale}{HTML}{e6f5f3}
\colorlet{colDef}{qteal}
\definecolor{colProp}{HTML}{0f766e}
\definecolor{colEx}{HTML}{b45309}
\definecolor{colRem}{HTML}{9d174d}
\color{ink}

\graphicspath{{../../../../assets/}}
\newcommand{\logofile}{../../../../assets/qimma-logo.pdf}

% ------------------- METADONNEES -------------------
\newcommand{\matiere}{Mathématiques}
\newcommand{\niveau}{2\up{ème} Bac -- Sciences Mathématiques}
\newcommand{\chapitretitre}{Fonctions primitives}

% ------------------- EN-TETE / PIED DE PAGE -------------------
\pagestyle{fancy}
\fancyhf{}
\renewcommand{\headrulewidth}{0pt}
\renewcommand{\footrulewidth}{0.5pt}
\fancyfoot[L]{\raisebox{-0.15cm}{\includegraphics[height=0.35cm]{\logofile}}~\small\sffamily\color{inkgray}Qimma}
\fancyfoot[C]{\small\sffamily\color{inkgray}\matiere\ -- \niveau}
\fancyfoot[R]{\small\sffamily\color{qtealdark}\thepage}

% ------------------- TITRES DE SECTION -------------------
\titleformat{\section}
  {\normalfont\sffamily\Large\bfseries\color{qtealdark}}
  {\tikz[baseline=-3.2pt]{\node[circle,fill=qteal,text=white,inner sep=0pt,minimum size=8mm]{\sffamily\bfseries\thesection};}}
  {10pt}{}
  [\vspace{-2mm}{\color{qamber}\titlerule[1.6pt]}\vspace{2mm}]
\titlespacing*{\section}{0pt}{16pt}{10pt}

\titleformat{\subsection}
  {\normalfont\sffamily\large\bfseries\color{qtealdark}}
  {\color{qteal}\ding{212}}{6pt}{}
\titlespacing*{\subsection}{0pt}{12pt}{6pt}

% ------------------- BOITES PEDAGOGIQUES -------------------
\newtcolorbox{fiche}[3][]{
  enhanced, breakable,
  colback=white, colframe=#2,
  boxrule=1pt, arc=2mm,
  top=7mm, left=4mm, right=4mm, bottom=3mm,
  drop shadow={black!25, shadow xshift=1mm, shadow yshift=-1mm},
  attach boxed title to top left={xshift=4mm, yshift=-3mm, yshifttext=-1mm},
  boxed title style={colback=#2, arc=1.5mm, boxrule=0pt, left=2mm, right=2mm},
  coltitle=white, fonttitle=\sffamily\bfseries,
  title={#3}, #1
}
\newenvironment{definition}[1][Définition]
  {\begin{fiche}{colDef}{\ding{110}~~#1}}{\end{fiche}}
\newenvironment{propriete}[1][Propriété]
  {\begin{fiche}{colProp}{\ding{72}~~#1}}{\end{fiche}}
\newenvironment{exemple}[1][Exemple]
  {\begin{fiche}{colEx}{\ding{217}~~#1}}{\end{fiche}}
\newenvironment{remarque}[1][Remarque]
  {\begin{fiche}{colRem}{\ding{43}~~#1}}{\end{fiche}}

% Démonstration (hors boîte, style discret)
\newenvironment{demo}
  {\par\smallskip\noindent{\sffamily\bfseries\color{qtealdark}Démonstration.}\ \itshape}
  {\ \normalfont\hfill$\blacksquare$\par\medskip}

\hypersetup{colorlinks=true, linkcolor=qtealdark, urlcolor=qtealdark}

% ============================================================
\begin{document}

% ------------------- BANNIERE DE TITRE -------------------
\begin{tcolorbox}[
  enhanced, boxrule=0pt, arc=4mm,
  interior style={left color=qtealdark, right color=qteal},
  top=8mm, bottom=8mm, left=10mm, right=32mm,
  drop shadow={black!35, shadow xshift=1.5mm, shadow yshift=-1.5mm},
  before skip=0pt, after skip=9mm,
  overlay={
    \node[anchor=north west] at ([xshift=6mm,yshift=-6mm]frame.north west)
      {\includegraphics[height=1.25cm]{\logofile}};
    \node[anchor=north east, fill=qamber, text=qambertext, rounded corners=2pt,
          inner sep=4pt, font=\sffamily\bfseries\small] at ([xshift=-6mm,yshift=-6mm]frame.north east) {COURS};
  }
]
\hspace{1.6cm}\centering
{\sffamily\bfseries\color{white}\fontsize{23}{27}\selectfont \chapitretitre}\\[3mm]
{\sffamily\color{white!90}\large \matiere\ \textbullet\ \niveau}
\end{tcolorbox}

% ------------------- OBJECTIFS -------------------
\begin{tcolorbox}[
  enhanced, colback=tealpale, colframe=qteal, boxrule=0.8pt, arc=2mm,
  left=5mm, right=5mm, top=3mm, bottom=3mm,
  title={\sffamily\bfseries\color{qtealdark}\ding{51}~~Objectifs du chapitre},
  coltitle=qtealdark, colbacktitle=tealpale, boxrule=0.8pt,
  attach title to upper={\par\vspace{1mm}}
]
\begin{itemize}[leftmargin=6mm, label=\ding{212}, itemsep=1mm]
  \item Démontrer que deux primitives d'une même fonction sur un intervalle diffèrent d'une constante.
  \item Connaître parfaitement les primitives des fonctions usuelles et savoir les redémontrer par dérivation.
  \item Reconnaître et exploiter toutes les formes composées ($u'u^n$, $\frac{u'}{u}$, $u'e^u$, $\frac{u'}{\sqrt u}$, $u'\cos u$, $u'\sin u$, $\frac{u'}{1+u^2}$).
  \item Déterminer la primitive vérifiant une condition initiale donnée.
  \item Combiner ajustement par une constante et décomposition d'une fraction pour trouver une primitive non directement reconnaissable.
  \item Résoudre les exercices type examen national mêlant primitives et fonctions ln/exp/trigonométriques.
\end{itemize}
\end{tcolorbox}

\vspace{4mm}

% ------------------- SECTION 1 -------------------
\section{Notion de primitive}

\begin{definition}
Soit $f$ une fonction définie sur un intervalle $I$. Une fonction $F$, dérivable sur $I$, est une \textbf{primitive de $f$ sur $I$} lorsque
\[
  F'(x) = f(x) \qquad \text{pour tout } x \in I.
\]
\end{definition}

\begin{propriete}[Existence]
Toute fonction \textbf{continue} sur un intervalle $I$ admet des primitives sur $I$ (théorème admis en terminale ; il sera justifié en analyse plus poussée via le calcul intégral).
\end{propriete}

\begin{propriete}[Structure de l'ensemble des primitives]
Si $F$ est une primitive de $f$ sur $I$, alors les primitives de $f$ sur $I$ sont exactement les fonctions $x\mapsto F(x)+C$, $C\in\mathbb{R}$.
\end{propriete}

\begin{demo}
\emph{Sens direct} : si $G(x)=F(x)+C$, alors $G'(x)=F'(x)=f(x)$ : $G$ est bien une primitive de $f$.

\smallskip
\emph{Réciproque} : soit $G$ une autre primitive de $f$ sur $I$. Posons $\varphi = G - F$. Alors $\varphi'(x) = G'(x)-F'(x) = f(x)-f(x) = 0$ pour tout $x\in I$. Comme $\varphi'=0$ sur l'\textbf{intervalle} $I$, $\varphi$ est \textbf{constante} sur $I$ (conséquence du théorème des accroissements finis, vu au chapitre sur la dérivation) : il existe $C\in\mathbb{R}$ tel que $\varphi(x)=C$ pour tout $x$, soit $G(x)=F(x)+C$.
\end{demo}

\begin{propriete}[Condition initiale]
Pour tout $(x_0,y_0)$ avec $x_0\in I$, il existe une \textbf{unique} primitive $F$ de $f$ sur $I$ telle que $F(x_0)=y_0$.
\end{propriete}

\begin{demo}
Soit $F_0$ une primitive quelconque de $f$ sur $I$. Toute primitive s'écrit $F=F_0+C$. La condition $F(x_0)=y_0$ donne $C=y_0-F_0(x_0)$, valeur \textbf{unique} : la primitive cherchée est $F=F_0+\bigl(y_0-F_0(x_0)\bigr)$.
\end{demo}

\begin{exemple}[Primitive avec condition initiale]
Déterminons la primitive $F$ de $f(x)=3x^2-4x+1$ sur $\mathbb{R}$ telle que $F(1)=2$.

\smallskip
Les primitives de $f$ sont $F(x)=x^3-2x^2+x+C$. La condition $F(1)=2$ donne $1-2+1+C=2\iff C=2$.

\smallskip
\textbf{Conclusion} : $F(x)=x^3-2x^2+x+2$.
\end{exemple}

\begin{remarque}[La condition « sur un intervalle » est essentielle]
Une primitive se définit toujours \textbf{sur un intervalle}. Sur $\mathbb{R}^*$ (réunion de deux intervalles, pas un intervalle), la « constante » de l'énoncé peut être \textbf{différente} sur $]-\infty\,;0[$ et sur $]0\,;+\infty[$ : par exemple, les primitives de $\frac1x$ sur $\mathbb{R}^*$ sont $\ln|x|+C_1$ sur $]-\infty\,;0[$ et $\ln|x|+C_2$ sur $]0\,;+\infty[$, avec $C_1$ et $C_2$ \textbf{indépendantes}.
\end{remarque}

% ------------------- SECTION 2 -------------------
\section{Primitives des fonctions usuelles}

\begin{propriete}[Tableau des primitives usuelles]
\begin{center}
\renewcommand{\arraystretch}{1.55}
\sffamily\small
\begin{tabular}{>{\centering\arraybackslash}p{4.5cm} >{\centering\arraybackslash}p{5.2cm} >{\centering\arraybackslash}p{5.2cm}}
\rowcolor{qtealdark} \color{white}\bfseries Fonction $f$ & \color{white}\bfseries Primitives $F$ ($C\in\mathbb{R}$) & \color{white}\bfseries Sur l'intervalle \\
\rowcolor{tealpale} $a$ (constante) & $ax+C$ & $\mathbb{R}$ \\
$x^n$ ($n\in\mathbb{N}$) & $\dfrac{x^{n+1}}{n+1}+C$ & $\mathbb{R}$ \\
\rowcolor{tealpale} $\dfrac{1}{x^n}$ ($n\geq2$) & $\dfrac{-1}{(n-1)x^{n-1}}+C$ & $]0\,;+\infty[$ ou $]-\infty\,;0[$ \\
$\dfrac{1}{\sqrt x}$ & $2\sqrt x+C$ & $]0\,;+\infty[$ \\
\rowcolor{tealpale} $\dfrac1x$ & $\ln|x|+C$ & $]0\,;+\infty[$ ou $]-\infty\,;0[$ \\
$e^x$ & $e^x+C$ & $\mathbb{R}$ \\
\rowcolor{tealpale} $\cos x$ & $\sin x+C$ & $\mathbb{R}$ \\
$\sin x$ & $-\cos x+C$ & $\mathbb{R}$ \\
\rowcolor{tealpale} $1+\tan^2x=\dfrac{1}{\cos^2x}$ & $\tan x+C$ & $\left]-\frac\pi2\,;\frac\pi2\right[$ \\
$\dfrac{1}{1+x^2}$ & $\arctan x+C$ & $\mathbb{R}$ \\
\rowcolor{tealpale} $\cos(ax+b)$ ($a\ne0$) & $\dfrac1a\sin(ax+b)+C$ & $\mathbb{R}$ \\
$\sin(ax+b)$ ($a\ne0$) & $-\dfrac1a\cos(ax+b)+C$ & $\mathbb{R}$ \\
\rowcolor{tealpale} $e^{ax+b}$ ($a\ne0$) & $\dfrac1ae^{ax+b}+C$ & $\mathbb{R}$ \\
\end{tabular}
\end{center}
\end{propriete}

\begin{remarque}[$\arctan$ hors-programme mais utile]
La primitive $\arctan x$ de $\frac{1}{1+x^2}$ figure ici pour information (fonction Arctan, notion approfondie en CPGE) ; au niveau bac SM, cette forme apparaît surtout via la composée $\frac{u'}{1+u^2}$ dont la primitive est $\arctan u+C$, rarement exigée mais parfois rencontrée en exercice guidé.
\end{remarque}

% ------------------- SECTION 3 -------------------
\section{Primitives des formes composées}

\begin{propriete}[Formes composées à reconnaître]
Si $u$ est dérivable sur $I$ (avec conditions de signe indiquées) :
\begin{center}
\renewcommand{\arraystretch}{1.55}
\sffamily\small
\begin{tabular}{>{\centering\arraybackslash}p{4.8cm} >{\centering\arraybackslash}p{5.6cm} >{\centering\arraybackslash}p{4.2cm}}
\rowcolor{qtealdark} \color{white}\bfseries Forme & \color{white}\bfseries Primitives & \color{white}\bfseries Condition \\
\rowcolor{tealpale} $u'u^n$ ($n\in\mathbb{N}$) & $\dfrac{u^{n+1}}{n+1}+C$ & -- \\
$\dfrac{u'}{u^n}$ ($n\geq2$) & $\dfrac{-1}{(n-1)u^{n-1}}+C$ & $u\ne0$ \\
\rowcolor{tealpale} $\dfrac{u'}{\sqrt u}$ & $2\sqrt u+C$ & $u>0$ \\
$\dfrac{u'}{u}$ & $\ln|u|+C$ & $u\ne0$ \\
\rowcolor{tealpale} $u'e^{u}$ & $e^{u}+C$ & -- \\
$u'\cos u$ & $\sin u+C$ & -- \\
\rowcolor{tealpale} $u'\sin u$ & $-\cos u+C$ & -- \\
$\dfrac{u'}{1+u^2}$ & $\arctan u+C$ & -- \\
\end{tabular}
\end{center}
\end{propriete}

\begin{demo}
Chaque formule se vérifie par dérivation directe. Par exemple pour $\dfrac{u^{n+1}}{n+1}$ : en dérivant comme composée, $\left(\dfrac{u^{n+1}}{n+1}\right)' = \dfrac{(n+1)u'u^n}{n+1} = u'u^n$, ce qui confirme la formule. Même démarche pour les autres lignes.
\end{demo}

\begin{exemple}[Reconnaître la forme $u'u^n$]
Déterminons les primitives de $f(x)=x(x^2+1)^3$ sur $\mathbb{R}$.

\smallskip
Posons $u(x)=x^2+1$, $u'(x)=2x$. On écrit $f(x)=\frac12\times2x(x^2+1)^3=\frac12u'u^3$. D'où
\[
  F(x) = \frac12\cdot\frac{(x^2+1)^4}{4}+C = \frac{(x^2+1)^4}{8}+C.
\]
\emph{Vérification} : $F'(x)=\dfrac{4\times2x(x^2+1)^3}{8}=x(x^2+1)^3=f(x)$. $\checkmark$
\end{exemple}

\begin{exemple}[Reconnaître la forme $u'/u$]
Déterminons les primitives de $g(x)=\dfrac{2x+1}{x^2+x+3}$ sur $\mathbb{R}$.

\smallskip
Le trinôme $u(x)=x^2+x+3$ a pour discriminant $\Delta=1-12<0$ : toujours strictement positif. Comme $u'(x)=2x+1$, on reconnaît exactement $\frac{u'}{u}$ :
\[
  G(x)=\ln\left(x^2+x+3\right)+C.
\]
\end{exemple}

\begin{exemple}[Coefficient ajusté]
Déterminons les primitives de $f(x)=\dfrac{x}{x^2+1}$ sur $\mathbb{R}$.

\smallskip
On écrit $f=\frac12\cdot\frac{2x}{x^2+1}=\frac12\cdot\frac{u'}{u}$ avec $u=x^2+1$. D'où $F(x)=\frac12\ln(x^2+1)+C$.
\end{exemple}

\begin{remarque}[La règle d'or de l'ajustement]
On n'ajuste \textbf{jamais} par une fonction de $x$, uniquement par une \textbf{constante multiplicative}. Si le coefficient devant $u'$ n'est pas exactement $1$, on factorise la constante manquante en dehors de l'intégrale/primitive.
\end{remarque}

% ------------------- SECTION 4 -------------------
\section{Techniques combinées : décomposition et cas non directs}

\begin{propriete}[Décomposition d'une fraction rationnelle]
Quand une fraction ne se ramène pas directement à une forme du tableau, on peut la \textbf{décomposer} en somme de termes plus simples (division euclidienne, ou décomposition en éléments simples pour un dénominateur factorisé).
\end{propriete}

\begin{exemple}[Division puis primitive]
Déterminons les primitives de $f(x)=\dfrac{x^2}{x+1}$ sur $]-1\,;+\infty[$.

\smallskip
Par division euclidienne : $x^2 = (x+1)(x-1)+1$, donc $f(x)=x-1+\dfrac{1}{x+1}$.

\smallskip
\[
  F(x) = \frac{x^2}{2}-x+\ln(x+1)+C \qquad\text{(sur }]-1\,;+\infty[\text{, où }x+1>0\text{)}.
\]
\end{exemple}

\begin{exemple}[Décomposition en éléments simples]
Déterminons les primitives de $f(x)=\dfrac{1}{x(x+1)}$ sur $]0\,;+\infty[$.

\smallskip
On cherche $a,b$ tels que $\dfrac{1}{x(x+1)}=\dfrac ax+\dfrac{b}{x+1}$. En réduisant au même dénominateur : $1=a(x+1)+bx$. Pour $x=0$ : $a=1$. Pour $x=-1$ : $b=-1$. Donc $f(x)=\dfrac1x-\dfrac1{x+1}$.

\smallskip
\[
  F(x) = \ln x - \ln(x+1)+C = \ln\frac{x}{x+1}+C.
\]
\end{exemple}

\begin{exemple}[Linéarisation trigonométrique]
Déterminons les primitives de $f(x)=\cos^2 x$ sur $\mathbb{R}$.

\smallskip
On linéarise avec $\cos^2x=\dfrac{1+\cos(2x)}{2}$ :
\[
  F(x) = \frac{x}{2}+\frac{\sin(2x)}{4}+C.
\]
\end{exemple}

% ------------------- SECTION 5 -------------------
\section{Exercices corrigés -- niveau examen national SM}

\begin{exemple}[Exercice 1 -- primitive avec racine et condition initiale]
Déterminer la primitive $F$ de $f(x)=\dfrac{x}{\sqrt{x^2+1}}$ sur $\mathbb{R}$ vérifiant $F(0)=3$.

\smallskip
\textbf{Corrigé.} Posons $u=x^2+1>0$, $u'=2x$. $f=\frac12\cdot\frac{u'}{\sqrt u}$, dont une primitive est $\frac12\times2\sqrt u=\sqrt u$. Donc les primitives de $f$ sont $F(x)=\sqrt{x^2+1}+C$.

\smallskip
$F(0)=\sqrt1+C=1+C=3\iff C=2$. \textbf{Conclusion} : $F(x)=\sqrt{x^2+1}+2$.
\end{exemple}

\begin{exemple}[Exercice 2 -- primitive avec exponentielle composée]
Déterminer les primitives de $f(x)=(2x-1)e^{x^2-x}$ sur $\mathbb{R}$.

\smallskip
\textbf{Corrigé.} Posons $u=x^2-x$, $u'=2x-1$ : on reconnaît exactement $u'e^u$. Donc $F(x)=e^{x^2-x}+C$.
\end{exemple}

\begin{exemple}[Exercice 3 -- fonction trigonométrique composée]
Déterminer les primitives de $f(x)=\sin x\cos^3x$ sur $\mathbb{R}$.

\smallskip
\textbf{Corrigé.} Posons $u=\cos x$, $u'=-\sin x$. On écrit $f=-(-\sin x)\cos^3x=-u'u^3$. Donc une primitive de $u'u^3$ est $\frac{u^4}{4}$, d'où
\[
  F(x) = -\frac{\cos^4x}{4}+C.
\]
\end{exemple}

\begin{exemple}[Exercice 4 -- fraction avec numérateur de degré supérieur]
Déterminer les primitives de $f(x)=\dfrac{x^3}{x^2+1}$ sur $\mathbb{R}$.

\smallskip
\textbf{Corrigé.} Division : $x^3 = x(x^2+1)-x$, donc $f(x)=x-\dfrac{x}{x^2+1}$.
\[
  F(x) = \frac{x^2}{2}-\frac12\ln(x^2+1)+C.
\]
\end{exemple}

\begin{exemple}[Exercice 5 -- combinaison $\ln$ et rationnelle]
Déterminer la primitive $G$ de $g(x)=\dfrac{\ln x}{x}$ sur $]0\,;+\infty[$ vérifiant $G(1)=0$.

\smallskip
\textbf{Corrigé.} On reconnaît $g=u'u$ avec $u=\ln x$ (car $u'=\frac1x$). Une primitive de $u'u$ est $\dfrac{u^2}{2}$, donc $G(x)=\dfrac{(\ln x)^2}{2}+C$.

\smallskip
$G(1)=\dfrac{0}{2}+C=C=0$. \textbf{Conclusion} : $G(x)=\dfrac{(\ln x)^2}{2}$.
\end{exemple}

% ------------------- SECTION 6 -------------------
\section{Méthodes et pièges : la boîte à outils de l'examen}

\subsection{Ce qu'on te demande, ce que tu fais}

\begin{center}
\renewcommand{\arraystretch}{1.45}
\sffamily\small
\begin{tabular}{>{\raggedright\arraybackslash}p{7.2cm} >{\raggedright\arraybackslash}p{8cm}}
\rowcolor{qtealdark} \color{white}\bfseries Ce qu'on te demande & \color{white}\bfseries Ton réflexe \\
\rowcolor{tealpale} Reconnaître une forme composée & Identifier $u$, calculer $u'$, comparer au numérateur/facteur. \\
Le coefficient devant $u'$ n'est pas $1$ & Factoriser/ajuster par une \textbf{constante} uniquement, jamais par du $x$. \\
\rowcolor{tealpale} Fraction rationnelle non reconnaissable & Division euclidienne si degré num $\geq$ degré dénom, sinon décomposition en éléments simples. \\
Fonction trigonométrique type $\sin x\cos^nx$ & Poser $u=\cos x$ (ou $\sin x$), retrouver $u'u^n$. \\
\rowcolor{tealpale} Primitive avec condition initiale & Trouver la famille de primitives, puis résoudre $F(x_0)=y_0$ pour $C$. \\
Vérifier une primitive trouvée & Toujours redériver $F$ et comparer à $f$. \\
\end{tabular}
\end{center}

\subsection{Les pièges qui coûtent des points}

\begin{remarque}[Erreurs relevées chaque année par les correcteurs]
\begin{itemize}[leftmargin=6mm, itemsep=0.5mm]
  \item Oublier la valeur absolue dans $\ln|u|+C$ (nécessaire dès que $u$ peut être négatif sur l'intervalle considéré).
  \item Ajuster un coefficient par une expression dépendant de $x$ au lieu d'une constante : c'est \textbf{interdit}, la primitive obtenue serait fausse.
  \item Oublier la constante $+C$ dans la réponse finale (sauf si une condition initiale la détermine).
  \item Confondre primitive et dérivée : dériver au lieu d'intégrer, ou inversement, lors de la vérification.
  \item Ne pas préciser l'intervalle de validité de la primitive (essentiel dès qu'il y a un $\ln$ ou une racine, où le domaine est restreint).
  \item Oublier de vérifier le signe de $u$ avant d'utiliser $\frac{u'}{\sqrt u}\to 2\sqrt u$ ($u$ doit être $>0$).
\end{itemize}
\end{remarque}

% ------------------- RESUME -------------------
\vspace{4mm}
\begin{tcolorbox}[
  enhanced, boxrule=0pt, arc=2mm,
  interior style={left color=qtealdark, right color=qteal},
  left=6mm, right=6mm, top=4mm, bottom=4mm,
  fontupper=\color{white}\sffamily,
  title={\sffamily\bfseries\color{white}\ding{72}~~À retenir},
  colbacktitle=qtealdark, coltitle=white, boxrule=0pt
]
\begin{itemize}[leftmargin=6mm, label={\color{qamber}\ding{212}}, itemsep=1mm]
  \item $F$ primitive de $f$ sur $I$ $\iff F'=f$ sur $I$ ; toute fonction continue admet des primitives ; deux primitives diffèrent d'une constante (démontré via TAF).
  \item Condition initiale $F(x_0)=y_0$ fixe la constante de façon unique.
  \item Formes clés : $u'u^n\to\frac{u^{n+1}}{n+1}$ ; $\frac{u'}{u}\to\ln|u|$ ; $u'e^u\to e^u$ ; $\frac{u'}{\sqrt u}\to2\sqrt u$ ; $u'\cos u\to\sin u$ ; $u'\sin u\to-\cos u$.
  \item Méthode : identifier $u$, calculer $u'$, ajuster par une \textbf{constante} uniquement, puis vérifier en dérivant.
  \item Fraction non reconnaissable : division euclidienne (degré num $\geq$ dénom) ou décomposition en éléments simples.
\end{itemize}
\end{tcolorbox}

\end{document}
